\chapter{Additional Experiments}

\section{Preliminary Exploration of Compression Techniques}
\label{app:preliminary-compression}

Before coming to the final decision on gradient accumulation with percentile-based sparsification, we explored a few other directions where communication might be reduced without losing too much gradient information. Both of these approaches seek to compress gradients based on the observation that gradients from different samples are not all equally informative.

\textbf{Summary Statistics Compression.} Here we replaced the full $(b\times 4096)$ gradient tensor sent from Machine~1 with compact summary statistics. Instead of providing the entire gradient tensor, we computed the mean and standard deviation across the batch dimension, producing a $(2\times 4096)$ summary. Machine~0 then reconstructed an approximate gradient tensor by sampling from $\mathcal{N}(\mu_j, \sigma_j)$ per dimension $j$. Sending only $2\times 4096$ statistics instead of the full $b \times 4096$ tensor gives us a compression factor of $b/2$ (64$\times$ at batch size 128). However, the reconstructed gradients are stochastic approximations that discard all per-sample structure, replacing the actual gradient landscape with noise shaped by the summary statistics.

\textbf{Information-Theoretic Filtering.} Our next idea built off of the summary statistics approach by maintaining per-dimension running mean $\mu_j$ and variance $\sigma_j^2$ over all minibatches using an exponential moving average. Instead of sending the summary statistics directly, we used these running distributions to filter individual gradient values based on how surprising they are. For each element in the full gradient tensor, we computed a z-score $z_{ij} = |g_{ij} - \mu_j| / \sigma_j$ measuring how unexpected that gradient value is relative to the running distribution. Only the top 10\% of elements by z-score were transmitted as sparse tensors. For unsent positions, Machine~0 reconstructed values by sampling from $\mathcal{N}(\mu_j, \sigma_j)$ rather than replacing them with zeros, providing a distributional approximation for the missing values. The running statistics were periodically synced from Machine~1 to Machine~0 at minimal overhead, requiring only two vectors of length 4096.

Figure~\ref{fig:early-experiments} shows training performance for both approaches against the uncompressed baseline over approximately 500,000 steps. The summary statistics compression comes in at roughly half the baseline, reaching an EMA return of 254 versus the baseline's 512. This indicates that the stochastic reconstruction preserves enough gradient signal for partial learning but loses too much per-sample information for competitive performance. The information-theoretic filtering approach fails to learn meaningfully, with episodic returns remaining near zero throughout training. Selecting gradients by how unusual they are relative to the running distribution does not provide a useful learning signal for the CNN feature extractor in this setting.

These results motivated our transition to gradient accumulation with percentile-based sparsification, which operates on raw gradient magnitudes and avoids stochastic reconstruction entirely.

\begin{figure}[htbp]
\centerline{\includegraphics[width=\textwidth]{images/exp5_early_experiments.pdf}}
\caption{Preliminary compression approaches over approximately 500,000 steps. Summary statistics compression achieves partial learning at roughly half the baseline return. Information-theoretic filtering fails to learn. The dashed vertical line indicates warm start completion at 30,000 steps.}
\label{fig:early-experiments}
\end{figure}

\section{Percentile Threshold Sweep}
\label{app:percentile-sweep}

In Figure~\ref{fig:exp6-returns}, we have shown results from training performance over five configurations using a random seed other than the one used in the primary experiments. We can see that the baseline and the 90th percentile with 0.99 decay are both able to converge, which confirms that the decay variant is similarly robust to the baseline. However, the higher threshold percentiles (the 95th percentile and the 99th percentile) with decay display increasingly poor performance. Therefore, it appears that transmitting less than ten percent of the gradient information does not allow for stable learning in this environment.

\begin{figure}[htbp]
\centerline{\includegraphics[width=\textwidth]{images/exp6_episodic_returns.pdf}}
\caption{Training performance for various percentile thresholds using a different random seed than the main experiments. The 90th percentile with 0.99 decay has similar convergence characteristics as the baseline. In contrast, we observe slow learning for the 95th percentile, while no learning occurs at all for the 99th percentile. The dashed vertical line indicates warm start completion at 30,000 steps.}
\label{fig:exp6-returns}
\end{figure}
